\documentclass[11pt,a4paper]{article}

\usepackage{fontspec}
\setmainfont{Times New Roman}
\setsansfont{Arial}
\setmonofont{Menlo}
\usepackage[russian]{babel}
\usepackage{indentfirst}
\usepackage{geometry}
\geometry{left=18mm,right=18mm,top=15mm,bottom=16mm}
\usepackage{microtype}
\usepackage{amsmath,amssymb}
\usepackage{enumitem}
\usepackage{hyperref}
\usepackage{caption}
\usepackage{graphicx}
\usepackage{xcolor}
\usepackage{tikz}
\usepackage{titlesec}
\usetikzlibrary{arrows.meta,positioning,calc}

\hypersetup{
  colorlinks=true,
  linkcolor=black,
  citecolor=black,
  urlcolor=black,
  pdftitle={Персонализированное ранжирование CrossSell-предложений}
}
\setlength{\parindent}{1.0em}
\setlength{\parskip}{0.12em}
\setlength{\emergencystretch}{2.5em}
\setlist[itemize]{leftmargin=1.20em,itemsep=0.04em,topsep=0.08em}
\captionsetup{font=small,labelfont=bf}
\setlength{\textfloatsep}{0.45em}
\setlength{\intextsep}{0.40em}
\setlength{\abovecaptionskip}{0.22em}
\setlength{\belowcaptionskip}{0.08em}
\titleformat{\section}{\large\bfseries}{\thesection}{0.65em}{}
\titlespacing*{\section}{0pt}{1.0ex plus .2ex}{0.45ex}

\newcommand{\NPV}{\mathrm{NPV}}
\newcommand{\E}{\mathbb{E}}
\newcommand{\I}{\mathbf{1}}
\newcommand{\AUC}{\mathrm{AUC}}
\newcommand{\sigm}{\sigma}
\newcommand{\Tr}{\operatorname{Tr}}

\tikzset{
  modelbox/.style={
    draw=#1!58!black,
    fill=#1!7,
    rounded corners=1.7mm,
    line width=0.65pt,
    align=center,
    minimum height=8.5mm,
    inner xsep=3.5pt,
    inner ysep=3.0pt,
    font=\scriptsize
  },
  modelboxsmall/.style={
    draw=#1!58!black,
    fill=#1!6,
    rounded corners=1.5mm,
    line width=0.60pt,
    align=center,
    minimum height=6.8mm,
    inner xsep=3.0pt,
    inner ysep=2.4pt,
    font=\scriptsize
  },
  flow/.style={-{Latex[length=2.1mm]}, line width=0.7pt, draw=black!58},
  auxflow/.style={-{Latex[length=1.8mm]}, line width=0.55pt, draw=black!48}
  ,
  pipeframe/.style={
    draw=#1!45!black,
    fill=#1!4,
    rounded corners=2.2mm,
    line width=0.70pt
  },
  pipebox/.style={
    draw=#1!55!black,
    fill=#1!7,
    rounded corners=1.8mm,
    line width=0.65pt,
    align=center,
    minimum height=9.5mm,
    inner xsep=4.0pt,
    inner ysep=3.0pt,
    font=\scriptsize
  },
  pipesmall/.style={
    draw=#1!55!black,
    fill=#1!8,
    rounded corners=1.5mm,
    line width=0.60pt,
    align=center,
    minimum height=8.0mm,
    inner xsep=3.5pt,
    inner ysep=2.6pt,
    font=\scriptsize
  }
}

\begin{document}
\linespread{0.965}\selectfont
\raggedbottom

\begin{center}
{\bfseries Не совпадает со статьей на конкурс JMLC}\\[1.0mm]
{\Large\bfseries Персонализированное ранжирование CrossSell-предложений}\\[1.2mm]
{\large Выбор категорий по extra NPV}\\[1.8mm]
Проект для Junior ML Contest 2026
\end{center}

\section{Бизнес-постановка}

В проекте решается задача персонализации CrossSell-раздачи в крупной финтех-экосистеме. Пользователю показывается не весь каталог, а несколько допустимых кросс-продуктовых предложений: сервисы, подписки, финансовые продукты или партнерские категории. Внутренние названия продуктов, точные значения ценности, схемы таблиц и рабочий код не раскрываются из-за NDA; в открытой версии используются обобщенные категории и синтетические данные. Бизнес-ценность подхода состоит в росте утилизации продуктов, удержании пользователя внутри экосистемы и более эффективном использовании бюджета скидок.

Цель раздачи не клик и не сам факт выбора категории, а \textit{утилизация}: пользователь должен начать пользоваться новым продуктом, счетом, сервисом или подпиской. Для пользователя $i$ известны признаки $x_i$, допустимые категории $K_i$, число слотов $m$ и нормированная ценность утилизации $\NPV_k$. Стратегия выбирает ровно $m$ категорий с учетом продуктовых правил: категория должна быть доступна, а две категории одной продуктовой группы не должны попасть в одну раздачу.

Первоначально ценность категории оценивалась как вероятность утилизации, умноженная на $\NPV_k$. Такая модель отвечает на вопрос: ``что пользователь, вероятно, утилизирует?'' Однако клиент может иметь высокую вероятность утилизации и без показа. Поэтому для выбора категорий используется оценка дополнительной ценности, созданной именно воздействием:
\[
extra\_NPV_{ik}=
\NPV_k\left[
\hat p(u_{ik}=1\mid x_i,k,z=1)-
\hat p(u_{ik}=1\mid x_i,k,z=0)
\right],
\]
где $z=1$ означает показ категории, а $z=0$ означает отсутствие показа. Разность вероятностей является оценкой uplift, то есть дополнительного эффекта показа.

\section{Данные и временные срезы}

В контуре используются два типа данных. Первый тип, табличная витрина для бустинга, содержит строки пользователь--категория--период, признаки до подготовки раздачи, доступность категории, назначение $z_{ik}$, вероятность назначения, факт утилизации и временной срез. Второй тип, исторические события для трансформера, представляет собой последовательность записей с временем события, типом действия, связанным объектом или категорией и дополнительными атрибутами. В нее попадают действия в приложении, взаимодействия с предложениями, транзакции, открытия и закрытия продуктов, связанные категории и сервисы.

Для каждого целевого месяца используется только история, доступная на момент подготовки раздачи. Последовательность для трансформера собирается из событий последних месяцев с учетом производственного лага; обучение, валидация и итоговая проверка разделяются по времени, а не случайным перемешиванием строк. Это сохраняет одинаковую схему обучения и применения модели и защищает признаки от утечки будущей утилизации.

\section{Модели: бустинг, трансформер и гибрид}

\textbf{Исходные модели вероятности утилизации.} В модельном контуре есть три уровня. Градиентный бустинг обучается на табличной витрине: агрегированных признаках пары пользователь--категория, частотах и давностях событий, истории взаимодействий, признаках категории и клиента. Он устойчив на табличных данных и дает сильную базовую точку сравнения.

Вторая модель, нейросетевая ранжирующая модель на базе трансформера, близка к two-tower архитектуре. Левая часть строит представление пользователя по последовательности событий: $h_i=\mathrm{BERT}_\theta(events_i)$. Правая часть строит представление действия: для обычного прогноза это категория, а для uplift это пара ``категория--показ'', то есть эмбеддинг действия содержит и категорию $k$, и индикатор $z$. Выходная MLP-голова оценивает вероятность утилизации:
\[
\hat p^{tr}_{ikz}=
\sigm\!\left(G_\theta\left[h_i,a_{kz},h_i\odot a_{kz},|h_i-a_{kz}|\right]\right).
\]
Такой трансформер учитывает порядок и контекст событий. Например, недавняя последовательность ``интерес к категории, действие в приложении, транзакция'' несет иной сигнал, чем те же события, разнесенные на месяцы.

\textbf{Uplift и гибридная модель.} Оценка uplift опирается на случайную раздачу, которая проводится каждый месяц на части аудитории внутри допустимого набора категорий и сохраняет те же ограничения, что и рабочая раздача. После запуска проверяются баланс предраздаточных признаков, доли назначений по категориям и группам, а также способность простой модели предсказывать факт показа по признакам до раздачи: $\AUC(\hat z(x_i))\approx 0.5$.

На логах такой раздачи обучается S-learner: модель получает признаки пары пользователь--категория и индикатор показа $z$, а таргетом остается факт утилизации. Гибридная модель не заменяет трансформер, а добавляет его прогноз как признак в итоговый бустинг вместе с табличными признаками и индикатором показа. Для одной пары пользователь--категория модель запускается при $z=1$ и $z=0$; разность прогнозов формирует uplift и extra NPV:
\[
\hat\tau_{ik}=\hat p^{hyb}_{ik,1}-\hat p^{hyb}_{ik,0},
\qquad
extra\_NPV_{ik}=\NPV_k\hat\tau_{ik}.
\]
Если назначение хорошо предсказывается по предраздаточным признакам, значит в логах осталась скрытая селекция и причинная интерпретация uplift ненадежна.

\textbf{Выбор категорий.} На применении модель скорирует пары пользователь--категория, после чего к результатам применяется преселект доступных категорий. Набор заполняется жадно: на каждом шаге выбирается категория с максимальной оценкой среди доступных, после чего исключаются категории, нарушающие правила сочетаний. Сравниваются случайный выбор, ранжирование по $\NPV$, ранжирование по ожидаемой ценности $\hat p(u_{ik}=1)\cdot \NPV_k$ и ранжирование по extra NPV.

\section{Проверка на исторических данных и открытая версия}

Стратегия выбора категорий по extra NPV проверяется на логах случайной раздачи. Основная практическая оценка считается на уровне выбранных категорий: если категория, выбранная новой стратегией, совпала с фактически показанной в случайной раздаче, ее вклад взвешивается через вероятность показа. Для более строгой проверки также считается оценка на уровне всего набора $A_i$, где $p_i(A_i)$ означает вероятность случайного назначения именно этого набора при тех же ограничениях. IPS-оценка ожидаемой ценности стратегии $\pi$ имеет вид:
\[
\widehat V_{\mathrm{IPS}}(\pi)=
\frac{1}{N}\sum_{i=1}^{N}
\frac{y_i\,\I\{A_i=\pi(x_i)\}}{p_i(A_i)}.
\]
Здесь $y_i$ обозначает наблюдаемый результат: утилизацию или нормированную ценность. Если раздача строится последовательно по слотам, вероятность набора восстанавливается из условных вероятностей назначений по слотам. Дополнительно считаются нормированная IPS-оценка, эффективный размер выборки, число совпадений стратегии с логом, бутстреп-интервалы и диагностика качества uplift.

Публичный репозиторий не содержит рабочих данных, точных NPV и внутреннего кода. Вместо этого в нем реализована синтетическая реконструкция: генератор пользователей, категорий, событий, случайной раздачи и неоднородного эффекта показа. В демо и notebook сравниваются случайный выбор, ранжирование по NPV, по $\hat p(u_{ik}=1)\cdot \NPV_k$, по extra NPV и выбор по истинному эффекту. В закрытом рабочем контуре последовательностный подход и extra NPV-ранжирование дают лучший порядок качества; открытая синтетика воспроизводит эту логику без раскрытия метрик и показывает, что модельные стратегии превосходят случайный выбор, а выбор по истинному эффекту задает верхнюю границу. Последняя стратегия возможна только в синтетике: она не является рабочим алгоритмом, а служит проверкой, что открытая версия учится выбирать добавочную ценность, а не воспроизводит общую склонность пользователя к активности.

\section{Инженерная часть}

Рабочее решение оформляется как параметризованный ML-пайплайн. Код пайплайна хранится в GitLab: отдельные Python-модули отвечают за загрузку данных, подготовку признаков, модели, оценку продуктов и формирование назначений. В CI код проходит проверки, после чего собирается Docker-образ для запуска расчета. Airflow оркестрирует сбор данных, расчет прогнозов для всех пар пользователь--категория и применение продуктовых правил.

\begin{center}
\begin{minipage}{\linewidth}
\centering
\resizebox{0.99\linewidth}{!}{%
\begin{tikzpicture}[x=1cm,y=1cm]
  \node[font=\bfseries\large] at (8.0,6.30) {ML-пайплайн};

  \node[pipebox=blue, text width=3.10cm] (gitlab) at (2.20,5.50)
    {GitLab\\[-0.5mm]{\tiny data, preprocessing, models, scoring, policy}};
  \node[pipebox=purple, text width=3.55cm] (ci) at (6.25,5.50)
    {CI проверяет код\\и собирает Docker-образ};
  \node[pipebox=orange, text width=4.15cm] (presel) at (11.55,5.50)
    {Преселекты от продуктов\\[-0.5mm]{\tiny доступные пары пользователь--категория}};

  \node[pipeframe=blue, minimum width=2.85cm, minimum height=3.15cm] (inputframe) at (1.60,2.58) {};
  \node[font=\bfseries\scriptsize] at (1.60,3.96) {Input};
  \node[pipesmall=blue, text width=2.20cm] (tabular) at (1.60,3.28)
    {Табличные данные\\по клиентам};
  \node[pipesmall=blue, text width=2.20cm] (events) at (1.60,2.08)
    {Данные по событиям\\клиентов};

  \node[pipeframe=teal, minimum width=10.55cm, minimum height=3.45cm] (airflowframe) at (8.50,2.58) {};
  \node[font=\bfseries\scriptsize, fill=teal!4, inner sep=1.0pt] at (8.50,3.98) {Сборка в Airflow};
  \node[pipebox=teal, text width=2.25cm] (data) at (4.75,2.65)
    {Сбор данных\\и обработка};
  \node[pipebox=green, text width=2.65cm] (scoring) at (8.15,2.65)
    {Скоринг всех пар\\пользователь--категория\\[-0.5mm]{\tiny boosting / transformer / stacking}};
  \node[pipebox=yellow, text width=3.05cm] (policy) at (11.85,2.65)
    {Оценка и выбор\\продуктов\\[-0.5mm]{\tiny score * NPV или extra NPV\\top-3 без повторов group\_id}};

  \node[pipeframe=gray, minimum width=2.35cm, minimum height=2.60cm] (outputframe) at (15.25,2.58) {};
  \node[font=\bfseries\scriptsize] at (15.25,3.67) {Output};
  \node[pipesmall=gray, text width=1.85cm] (output) at (15.25,2.48)
    {клиент\\номер слота\\продукт\\группа\\модель\\правило};

  \draw[flow] (gitlab) -- (ci);
  \draw[flow] (tabular.east) -- ([yshift=2.0mm]data.west);
  \draw[flow] (events.east) -- ([yshift=-2.0mm]data.west);
  \draw[flow] (data) -- (scoring);
  \draw[flow] (scoring) -- (policy);
  \draw[flow] (ci.south) -- ++(0,-0.48) -| (airflowframe.north);
  \draw[flow] (presel.south) -- ++(0,-0.35) -| (policy.north);
  \draw[flow] (airflowframe.east) -- (outputframe.west);
\end{tikzpicture}
}
\captionof{figure}{Инженерный контур расчета: код хранится в GitLab, CI собирает запускаемую версию, Airflow собирает данные, рассчитывает прогнозы и применяет правила выбора категорий.}
\end{minipage}
\end{center}

\section{Применение ИИ}

Для автоматизации регулярной аналитической работы я создал сценарий ИИ-аналитика. Он работает с результатами запусков, логами и построенным мной ML-дашбордом: ищет аномалии в метриках по моделям и продуктовым группам, формулирует гипотезы, готовит SQL-срезы, черновики задач и отчеты. Перед созданием задачи или отчета агент запрашивает подтверждение, а после выполнения кратко фиксирует, что сделал и какие проверки нужны дальше.

MCP-коннекторы и корпоративные доступы я не разрабатывал как отдельную платформу: моя роль состояла в интеграции доступных ИИ-инструментов в рабочий процесс и настройке правил их безопасного применения. Аналитический контур проще автоматизировать первым, потому что SQL-запросы и отчеты легче проверить, а цена ошибки ниже. Поэтому на аналитике обкатываются подходы, которые затем можно осторожно переносить в ML-разработку.

В ML-контуре уровень автоматизации ниже: изменение модели или стратегии ранжирования требует экспериментов, причинной проверки, контроля утечек и учета продуктовых ограничений. Здесь ИИ используется как набор специализированных навыков: проверка данных и временных срезов, генерация гипотез, ревью uplift-оценки, помощь с кодом и анализ ошибок в GitLab/CI. Для подготовки открытой версии проекта был создан отдельный навык, который проверяет техническую постановку, причинную интерпретацию, воспроизводимость кода и соблюдение NDA.

\section{Личный вклад и развитие}

Мой вклад включает формализацию CrossSell-раздачи как задачи ранжирования с ограничениями, разработку и сравнение бустинга и трансформерной модели, построение гибридной схемы со стекингом, внедрение оценки extra NPV для выбора категорий, проверку случайного логирования и создание открытой синтетической реконструкции с учетом NDA. Продуктовые правила, бизнес-метрики и запуск раздач обсуждаются совместно с менеджерами и аналитиками; модельная постановка, эксперименты и открытая версия находятся в моей зоне ответственности.

Следующий этап развития состоит в том, чтобы для части продуктов перейти от выбора только категории к совместному выбору категории и ставки скидки. Для каждой допустимой пары пользователь--категория можно сравнивать несколько уровней скидки и выбирать вариант, который максимизирует extra NPV за вычетом ожидаемой стоимости скидки с учетом общего бюджетного ограничения. Такой переход превращает ранжирование в задачу оптимизации действия ``категория--ставка скидки'', сохраняя те же правила доступности, разнообразия набора и проверку качества на случайных логах.

\end{document}
